\scanpage{85}
% manuscript line 1
\[
\mathop{\handwritten{\sum}}\nolimits_{j=0}
(x_1\,T_1+\ldots+x_{k-1}\,T_{k-1})^j\,x_k\,T_k\,
(x_1 T_1+\ldots+x_{k-1}\,T_{k-1})^{k-j-1}
\]

% manuscript line 2
Etant donnés deux éléments $x,y$ de $L$, on définit $g_{i,k-1}(x,y)$ comme le coefficient de $T_1^i\,T_2^{k-i}$ dans $(x\,T_1+y\,T_2)^k$. Montrer que l'on a
% manuscript line 4
\[
i!\,(k-i)!\,g_{i,k-i}(x,y)=f_k(t_1,\ldots,t_k)
\]
% manuscript line 5
avec $t_h=x$ pour $h\handwritten{\leqslant}i$ et $t_h=y$ pour $h\handwritten{\geqslant}i+1$.

% manuscript line 6
b) Lorsque $L$ est considérée comme algèbre de Lie sur $K$, on a , dans $\handwritten{\unsure{\mathscr L|\mathcal L}}_K(L)$\,:
% manuscript line 7
\[
(\mathrm{ad}\,x)^n=\mathop{\handwritten{\sum}}\nolimits_{i=0}^{n}
(-1)^{n-i}\binom{n}{i}\,L_x^i\,R_x^{n-i}.
\]
% manuscript line 8
En déduire que, si $K$ est de caractéristique $p$ ($p$ premier), on a\,:
% manuscript line 9
\begingroup
\makeatletter
\tagsleft@true
\makeatother
\[
(\mathrm{ad}\,x)^p.y=(\mathrm{ad}\,x^p).y
\tag*{(1)}
\]
% manuscript line 10
\[
(\mathrm{ad}\,x)^{p-1}.y=\mathop{\handwritten{\sum}}\nolimits_{i=0}^{p-1}x^i\,y\,x^{p-1-i}
\tag*{(2)}
\]
\endgroup

% manuscript line 11
Déduire de (2) que\,:
% manuscript line 12
\[
f_{p-1}(\mathrm{ad}\,x_1,\ldots,\mathrm{ad}\,x_{p-1}).y=f_p(x_1,\ldots,x_{p-1},y)
\]
% manuscript line 13
et\,:
% manuscript line 14
\[
g_{i,p-1-i}(\mathrm{ad}\,x,\mathrm{ad}\,y).y=(p-i)\,g_{i,p-i}(x,y).
\]
% archive sheet number: 84
